% notes/artefacts/market-making-lob/model-derivations.tex
% Full derivations behind four compendium entries, which state each result with its locator and an idea of
% the proof: mod:grossman-miller with res:grossman-miller-trading-fees and res:grossman-miller-liquidity-measures,
% mod:kyle-single-auction, mod:static-limit-order-depth and res:roll-spread-estimator. Notation follows the
% entries. The working is own derivation following the cited textbook accounts; the closed forms are checked
% numerically, with fixed parameters, seeds and tolerances, by model-derivations-checks.py in this folder
% (fee equilibrium by root-finding, the liquidity slope and autocovariance, Kyle's profit, projection and
% posterior variance by simulation, the limit-order argmax, and Roll's moments).
% Treat it as own working: source checking, not study.
% Compile inside notes/artefacts/market-making-lob/ with  latexmk -pdf model-derivations.tex
\PassOptionsToPackage{table}{xcolor}
\documentclass[a4paper,oneside,openany]{book}
% Your own packages, macros and environments. The learning guide's entry environments,
% \citep and \qfsource are in house.tex, which compendium.tex loads after this file.

% -----------------------------
% Document Class & Encoding
% -----------------------------
\usepackage[utf8]{inputenc}  % Input encoding
\usepackage[T1]{fontenc}     % Font encoding

% -----------------------------
% Page Layout & Geometry
% -----------------------------
\usepackage[margin=2cm,foot=1cm]{geometry}

% -----------------------------
% Header & Footer Management
% -----------------------------
\usepackage{fancyhdr}
\pagestyle{fancy}

% -----------------------------
% Text Formatting & Spacing
% -----------------------------
\usepackage{ragged2e}
\usepackage[nodisplayskipstretch]{setspace}
\setstretch{1}  % Line spacing

% -----------------------------
% Listings for Code
% -----------------------------
\usepackage{listings}
\usepackage{xcolor}

% -----------------------------
% Tables, Figures & Graphics
% -----------------------------
\usepackage{tabularx}
\usepackage{multirow}
\usepackage{pdflscape}
\usepackage{graphicx}

% -----------------------------
% Algorithms
% -----------------------------
\usepackage{algorithm}       % For floating algorithm environment
\usepackage{algpseudocode}   % For pseudo-code inside algorithms

% -----------------------------
% Math Packages
% -----------------------------
\usepackage{amsmath,amssymb,amsthm,mathtools,bm}
\usepackage{mathrsfs}        % Script fonts
\usepackage{tikz-cd}         % Commutative diagrams
\usepackage{rotating}        % Rotated objects
\usepackage{float}           % Float placement
\usetikzlibrary{calc, arrows, positioning, arrows.meta}

% -----------------------------
% TikZ and its Libraries
% -----------------------------
\usepackage{tikz}
\usetikzlibrary{positioning} % Added to support "below=of" syntax
\usepackage{forest}          % For advanced tree diagrams
\usepackage{tikz-qtree}      % For vertical tree diagrams

% -----------------------------
% Miscellaneous Packages
% -----------------------------
\usepackage{xparse}          % For defining robust commands
\usepackage{enumitem}        % For customizing list layouts
\usepackage{makecell}        % For enhanced table cells
\usepackage{comment}		 % For commenting out sections

%%%%%%%%%%%%%%%%%%%%%%%%%%%%%%%%%%%%%%%%%%%%%%%%%%%%%%%%%%%%%%%%%%%%%%%%%%%%%%%
% Custom Math Commands and Shortcuts
%%%%%%%%%%%%%%%%%%%%%%%%%%%%%%%%%%%%%%%%%%%%%%%%%%%%%%%%%%%%%%%%%%%%%%%%%%%%%%%

% Matrix and Vector Formatting
\newcommand{\matr}[1]{\mathbf{#1}}      % For undergraduate algebra (alternative: \bm{#1})
\newcommand{\vect}[1]{\bm{#1}}

% Standard Number Sets
\newcommand{\N}{\mathbb{N}}
\newcommand{\R}{\mathbb{R}}
\newcommand{\Q}{\mathbb{Q}}
\newcommand{\C}{\mathbb{C}}
\newcommand{\Z}{\mathbb{Z}}

% Bold variables (customize as needed)
\newcommand{\x}{\mathbf{x}}
\newcommand{\F}{\mathbf{F}}
\newcommand{\f}{\mathbf{f}}
\newcommand{\y}{\mathbf{y}}
\renewcommand{\a}{\mathbf{a}}
\renewcommand{\b}{\mathbf{b}}
\renewcommand{\c}{\mathbf{c}}
\newcommand{\h}{\mathbf{h}}
\newcommand{\g}{\mathbf{g}}
\newcommand{\z}{\mathbf{z}}
\newcommand{\ze}{\mathbf{0}}

% Norms, Absolute Values, and Bracketing
\newcommand{\norm}[1]{\left\lVert#1\right\rVert}
\newcommand{\abs}[1]{\left\lvert#1\right\rvert}
\newcommand{\brk}[1]{\left[#1\right]}
\newcommand{\brc}[1]{\left\{#1\right\}}
\newcommand{\paren}[1]{\left(#1\right)}
\newcommand{\normop}[1]{\left\lVert#1\right\rVert_\text{op}}

% Miscellaneous Math Shortcuts
\newcommand{\LL}{\mathcal{L}}
\newcommand{\uni}{\overset{\text{uni}}{\to}}
\DeclareMathOperator{\diam}{diam}
\newcommand{\Prr}[1]{\text{Pr}\left(#1\right)}
\newcommand{\ceil}[1]{\lceil #1 \rceil}
\newcommand{\floor}[1]{\lfloor #1 \rfloor}
\newcommand\disteq{\stackrel{\textnormal{dist}}{=}}

% Evaluation command with options
\NewDocumentCommand{\evalat}{sO{\big}mm}{%
  \IfBooleanTF{#1}
    {\mleft. #3 \mright|_{#4}}
    {#3#2|_{#4}}%
}

%%%%%%%%%%%%%%%%%%%%%%%%%%%%%%%%%%%%%%%%%%%%%%%%%%%%%%%%%%%%%%%%%%%%%%%%%%%%%%%
% Color Definitions
%%%%%%%%%%%%%%%%%%%%%%%%%%%%%%%%%%%%%%%%%%%%%%%%%%%%%%%%%%%%%%%%%%%%%%%%%%%%%%%

\definecolor{dartmouthgreen}{rgb}{0.05, 0.5, 0.06}
\definecolor{egyptianblue}{rgb}{0.06, 0.2, 0.65}
\definecolor{dukeblue}{rgb}{0.0, 0.0, 0.61}
\definecolor{jazzberryjam}{rgb}{0.65, 0.04, 0.37}
\definecolor{magenta}{HTML}{EC008C}
\definecolor{darkmagenta}{rgb}{0.55, 0.0, 0.55}
\definecolor{deeppink}{rgb}{1.0, 0.08, 0.58}
\definecolor{codegreen}{rgb}{0,0.6,0}
\definecolor{codegray}{rgb}{0.5,0.5,0.5}
\definecolor{codepurple}{rgb}{0.58,0,0.82}
\definecolor{backcolour}{rgb}{0.95,0.95,0.92}

%%%%%%%%%%%%%%%%%%%%%%%%%%%%%%%%%%%%%%%%%%%%%%%%%%%%%%%%%%%%%%%%%%%%%%%%%%%%%%%
% Custom Shortcuts and Emphasis Commands
%%%%%%%%%%%%%%%%%%%%%%%%%%%%%%%%%%%%%%%%%%%%%%%%%%%%%%%%%%%%%%%%%%%%%%%%%%%%%%%

\newcommand{\isomorp}{\xrightarrow{\sim}}
\newcommand{\hlt}[1]{\textit{{\color{blue}#1}}}
\newcommand{\simpt}[1]{\textit{{\color{deeppink}#1}}}
\newcommand{\impt}[1]{\textit{{\color{red}#1}}}
\newcommand{\gcds}[1]{\textnormal{gcd}#1}
\newcommand{\mins}[1]{\textnormal{min}#1}
\newcommand{\maxs}[1]{\textnormal{max}#1}
\newcommand{\lcms}[1]{\textnormal{lcm}#1}
\newcommand{\degs}[1]{\textnormal{deg}#1}
\newcommand{\exps}[1]{\textnormal{exp}#1}
\newcommand{\tors}[1]{\textnormal{Tor}#1}
\newcommand{\homs}[1]{\textnormal{Hom}#1}
\newcommand{\anns}[1]{\textnormal{Ann}#1}
\newcommand{\defeq}{\vcentcolon=}
\DeclareMathOperator{\sign}{sign}

%%%%%%%%%%%%%%%%%%%%%%%%%%%%%%%%%%%%%%%%%%%%%%%%%%%%%%%%%%%%%%%%%%%%%%%%%%%%%%%
% Theorem and Definition Environments
%%%%%%%%%%%%%%%%%%%%%%%%%%%%%%%%%%%%%%%%%%%%%%%%%%%%%%%%%%%%%%%%%%%%%%%%%%%%%%%
\theoremstyle{plain}
\newtheorem{theorem}{Theorem}[section]
\newtheorem{corollary}[theorem]{Corollary}
\newtheorem{lemma}[theorem]{Lemma}
\newtheorem{proposition}[theorem]{Proposition}
\newtheorem{qnbank}[theorem]{Question Bank}
\let\oldqnbank\qnbank
\renewcommand{\qnbank}{\oldqnbank\normalfont}
\newtheorem{process}[theorem]{Process}
\let\oldprocess\process
\renewcommand{\process}{\oldprocess\normalfont}
\newtheorem{method}[theorem]{Method}
\let\oldmethod\method
\renewcommand{\method}{\oldmethod\normalfont}
\newtheorem{definition}[theorem]{Definition}
\let\olddefinition\definition
\renewcommand{\definition}{\olddefinition\normalfont}
\newtheorem{example}[theorem]{Example}
\let\oldexample\example
\renewcommand{\example}{\oldexample\normalfont}
\newtheorem{remark}[theorem]{Remark}
\let\oldremark\remark
\renewcommand{\remark}{\oldremark\normalfont}
\newtheorem{notation}[theorem]{Notation}
\let\oldnotation\notation
\renewcommand{\notation}{\oldnotation\normalfont}
\newtheorem{principle}[theorem]{Principle}
\let\oldprinciple\principle
\renewcommand{\principle}{\oldprinciple\normalfont}
\newtheorem{mechanism}[theorem]{Mechanism}
\let\oldmechanism\mechanism
\renewcommand{\mechanism}{\oldmechanism\normalfont}
\newtheorem{taxonomy}[theorem]{Taxonomy}
\let\oldtaxonomy\taxonomy
\renewcommand{\taxonomy}{\oldtaxonomy\normalfont}
\newtheorem{model}[theorem]{Model}
\let\oldmodel\model
\renewcommand{\model}{\oldmodel\normalfont}
\newtheorem{assumption}[theorem]{Assumption}
\let\oldassumption\assumption
\renewcommand{\assumption}{\oldassumption\normalfont}
\newtheorem{casestudy}[theorem]{Case Study}
\let\oldcasestudy\casestudy
\renewcommand{\casestudy}{\oldcasestudy\normalfont}
\newtheorem{empiricalfinding}[theorem]{Empirical Finding}
\let\oldempiricalfinding\empiricalfinding
\renewcommand{\empiricalfinding}{\oldempiricalfinding\normalfont}
\newtheorem{practitionernote}[theorem]{Practitioner Note}
\let\oldpractitionernote\practitionernote
\renewcommand{\practitionernote}{\oldpractitionernote\normalfont}
\newtheorem{caveat}[theorem]{Caveat}
\let\oldcaveat\caveat
\renewcommand{\caveat}{\oldcaveat\normalfont}
\newtheorem{marketrule}[theorem]{Rule}
\let\oldmarketrule\marketrule
\renewcommand{\marketrule}{\oldmarketrule\normalfont}
\newtheorem{takeaway}[theorem]{Takeaway}
\let\oldtakeaway\takeaway
\renewcommand{\takeaway}{\oldtakeaway\normalfont}

%%%%%%%%%%%%%%%%%%%%%%%%%%%%%%%%%%%%%%%%%%%%%%%%%%%%%%%%%%%%%%%%%%%%%%%%%%%%%%%
% Breakable Algorithm Environment (Algorithms spanning multiple pages)
%%%%%%%%%%%%%%%%%%%%%%%%%%%%%%%%%%%%%%%%%%%%%%%%%%%%%%%%%%%%%%%%%%%%%%%%%%%%%%%

\makeatletter
\newenvironment{breakablealgorithm}
 {%
  \begin{center}
    \refstepcounter{algorithm} % Increment algorithm counter
    \hrule height.8pt depth0pt \kern2pt
    \renewcommand{\caption}[2][\relax]{%
      {\raggedright\textbf{Algorithm \thealgorithm:} ##2\par}%
      \ifx\relax##1\relax
        \addcontentsline{loa}{algorithm}{\protect\numberline{\thealgorithm}##2}%
      \else
        \addcontentsline{loa}{algorithm}{\protect\numberline{\thealgorithm}##1}%
      \fi
      \kern2pt\hrule\kern2pt
    }%
 }{%
    \kern2pt\hrule
  \end{center}
 }
\makeatother

%%%%%%%%%%%%%%%%%%%%%%%%%%%%%%%%%%%%%%%%%%%%%%%%%%%%%%%%%%%%%%%%%%%%%%%%%%%%%%%
% Code Listings Style
%%%%%%%%%%%%%%%%%%%%%%%%%%%%%%%%%%%%%%%%%%%%%%%%%%%%%%%%%%%%%%%%%%%%%%%%%%%%%%%

\lstdefinestyle{codestyle}{
  commentstyle=\color{codegreen},
  keywordstyle=\color{magenta},
  numberstyle=\tiny\color{codegray},
  stringstyle=\color{codepurple},
  basicstyle=\ttfamily\normalsize,
  breakatwhitespace=false,
  breaklines=true,
  captionpos=b,
  keepspaces=true,
  numbers=left,
  numbersep=10pt,
  frame=single,
  showspaces=false,
  showstringspaces=false,
  showtabs=false,
  tabsize=2
}
\lstset{style=codestyle}

% -----------------------------
% Hyperlinks (loaded last: loaded before float and algorithm it duplicated figure anchors)
% -----------------------------
\usepackage{hyperref}
\hypersetup{
  colorlinks,
  linkcolor={blue!80!black},
  citecolor={green!80!black},
  urlcolor={red!80!black}
}

\usepackage[T1]{fontenc}
\usepackage{amsmath,amssymb,booktabs,graphicx,hyperref}
\hypersetup{colorlinks}
\setlength{\emergencystretch}{3em}   % lets a line holding a long key break
\newcounter{qfentry}[chapter]
\renewcommand{\theqfentry}{\thechapter.\arabic{qfentry}}
% \qftype{type}{Name} declares the environment qf<type>; its title wraps like a paragraph.
\newcommand{\qftype}[2]{\newenvironment{qf#1}[1][]{\par\addvspace{\medskipamount}%
  \refstepcounter{qfentry}{\raggedright\noindent\bfseries #2 \theqfentry
  \if\relax\detokenize{##1}\relax\else\quad ##1\fi\par}\nopagebreak}%
  {\par\addvspace{\medskipamount}}}
\qftype{definition}{Definition}\qftype{convention}{Convention}\qftype{result}{Result}
\qftype{model}{Model}\qftype{method}{Method}\qftype{mechanism}{Mechanism}
\qftype{evidence}{Evidence}\qftype{data}{Data}\qftype{rule}{Rule}
\qftype{pitfall}{Pitfall}\qftype{example}{Example}
\newcommand{\qffield}[1]{\par\textbf{#1.}\enspace}   % starts one of an entry's fields
% \citep[locator]{key} links the key to its line in the sources if it has one,
% and otherwise to its entry in the pinned catalogue PDF.
\newcommand{\qfcatalogue}{../guide/institutional_quantitative_finance_reading_list.pdf}
\NewDocumentCommand{\citep}{o o m}{[\IfValueT{#2}{#1 }%
  \ifcsname r@src:#3\endcsname\hyperref[src:#3]{\nolinkurl{#3}}%
  \else\href{\qfcatalogue\##3}{\nolinkurl{#3}}\fi
  \IfValueTF{#2}{, #2}{\IfValueT{#1}{, #1}}]}
% \qfsource{key}{reference} records one source outside the catalogue, in front/sources.tex.
\newcommand{\qfsource}[2]{\par\noindent\phantomsection\label{src:#1}\nolinkurl{#1}\quad #2\par}
% \qfsections{key, ...} inputs <key>/section.tex for each listed key that has one, in order.
\ExplSyntaxOn
\NewDocumentCommand{\qfsections}{m}
  {\clist_map_inline:nn {#1} {\InputIfFileExists{##1/section}{}{}}}
\ExplSyntaxOff

% Catalogue links are relative to this folder, not to compendium/.
\renewcommand{\qfcatalogue}{../../../guide/institutional_quantitative_finance_reading_list.pdf}
\pagestyle{plain}
\begin{document}
\chapter*{Market making: model derivations}

\section*{Grossman--Miller market making}

Sources: \citep[chap.~``A Primer on the Microstructure of Financial Markets'', pp.~21--27]{crypto-cartea}; the free-entry condition from \citep[eqs.~(16)--(19), printed pp.~18--19]{grossman-miller-1988-nber}. Entries: \texttt{mod:grossman-miller}, \texttt{res:grossman-miller-trading-fees}, \texttt{res:grossman-miller-liquidity-measures}.

\subsection*{Assumptions}
\begin{enumerate}[label=\roman*.]
\setlength{\itemsep}{0pt}
\item Dates $t\in\{1,2,3\}$: date~1 is the first liquidity trade, date~2 the arrival of the offsetting liquidity trader, date~3 terminal cash settlement.
\item Cash earns nothing; the risky asset is worth $S_3=\mu+\epsilon_2+\epsilon_3$ at date~3, with $\epsilon_2,\epsilon_3$ independent $N(0,\sigma^2)$ and $\epsilon_t$ announced between dates $t-1$ and $t$.
\item $n$ identical market makers (MM) start with no inventory; liquidity trader LT1 holds $i$ units and LT2, present only from date~2, holds $-i$. For $i>0$ LT1 is a natural seller and LT2 a later buyer.
\item Every agent has CARA utility $U(X)=-\exp(-\gamma X)$, $\gamma>0$.
\item No trading, funding or holding costs, except where per-share fees are introduced below; the only cost of supplying liquidity is carrying inventory between dates~1 and~2. LT2's arrival at date~2 is certain.
\end{enumerate}
Let $q_t^j$ be agent $j$'s holding on leaving date $t$ and $S_t$ the date-$t$ trade price. Wealth evolves as
\[ X_2^j=X_1^j+q_1^j(S_2-S_1), \qquad X_3^j=X_2^j+q_2^j(S_3-S_2). \]

\subsection*{The CARA--normal demand}
If terminal wealth is $X+q(Z-S)$ with $Z\sim N(m,s^2)$ given the agent's information, then
\[ \mathbb{E}\big[U\big(X+q(Z-S)\big)\big]=-\exp\Big(-\gamma\big(X+q(m-S)\big)+\tfrac12\gamma^2q^2s^2\Big), \]
by the normal moment-generating function. Maximising expected utility is therefore maximising the certainty equivalent $X+q(m-S)-\tfrac12\gamma q^2s^2$, which is strictly concave in $q$; the first-order condition $m-S-\gamma qs^2=0$ gives the demand
\[ q^\star=\frac{m-S}{\gamma s^2}. \]

\subsection*{Date 2}
Given $\epsilon_2$, the next value is $S_3$ with $m=\mathbb{E}[S_3\mid\epsilon_2]=\mu+\epsilon_2$ and $s^2=\sigma^2$, so every agent present demands
\[ q_2^{j}=\frac{\mathbb{E}[S_3\mid\epsilon_2]-S_2}{\gamma\sigma^2}. \]
All $n+2$ agents are present. Holdings entering date~2 sum to $nq_1^{MM}+q_1^{LT1}+q_1^{LT2}=i-i=0$, because trading at date~1 only redistributes LT1's $i$ units and LT2 enters with $-i$. Market clearing requires the date-2 holdings to sum to the same total:
\[ 0=(n+2)\,\frac{\mathbb{E}[S_3\mid\epsilon_2]-S_2}{\gamma\sigma^2}
\quad\Longrightarrow\quad S_2=\mathbb{E}[S_3\mid\epsilon_2]=\mu+\epsilon_2,\qquad q_2^j=0. \]
Date~2 is the unwind: the matching liquidity has arrived and no agent wants risky inventory.

\subsection*{Date 1}
At date~1, LT2 is absent. Each present agent anticipates the date-2 price $S_2=\mu+\epsilon_2$, so $m=\mathbb{E}[S_2]=\mu$, $s^2=\operatorname{Var}(S_2)=\sigma^2$, and
\[ q_1^{j}=\frac{\mu-S_1}{\gamma\sigma^2},\qquad j\in\{MM,LT1\}. \]
Clearing among the $n$ market makers and LT1, whose holdings entering date~1 are $0$ and $i$:
\[ i=nq_1^{MM}+q_1^{LT1}=(n+1)\,\frac{\mu-S_1}{\gamma\sigma^2}, \]
so
\[ q_1^{MM}=q_1^{LT1}=\frac{i}{n+1},\qquad S_1=\mu-\gamma\sigma^2\,\frac{i}{n+1}. \]
LT1 sells $i-\frac{i}{n+1}=\frac{n}{n+1}\,i$ at date~1 and each market maker holds $\frac{i}{n+1}$ until LT2 arrives. The liquidity discount is $p=\abs{\mu-S_1}=\gamma\sigma^2\abs{i}/(n+1)$; for $i<0$ the same argument gives a premium. Interpretation: a bid at $\mu-p$ and an ask at $\mu+p$ imply an inside spread of $2p$.

The concession is a price for risk, not information: no agent is informed and the imbalance nets to zero across dates. It rises with $\abs{i}$, $\gamma$ and $\sigma^2$ and falls with $n$; as $n\to\infty$, $S_1\to\mu$ and LT1's immediate trade tends to $i$.

\emph{Free entry.} In the original, a market maker forgoes opportunities worth $c$, and entry occurs until $\mathbb{E}U\big(W_0-c+(\tilde P_2-P_1)x_1^m\big)=\mathbb{E}U(W_0)$ \citep[eq.~(16), printed p.~18]{grossman-miller-1988-nber}: an equality of expected utilities, not of expected spread capture and cost. Its consequences for the number of market makers are in \texttt{mod:grossman-miller}.

\subsection*{Per-share fees}
Let a fee $\eta$ be paid by each side per share traded. At date~2 a position is unwound either then or at date~3, so the fee enters both sides of each comparison and cancels:
\[ q_2^{j}=\frac{\mathbb{E}[S_3-\eta\mid\epsilon_2]-(S_2-\eta)}{\gamma\sigma^2}\ (j\in\{MM,LT1\}),\qquad
q_2^{LT2}=\frac{\mathbb{E}[S_3+\eta\mid\epsilon_2]-(S_2+\eta)}{\gamma\sigma^2}, \]
and clearing again gives $S_2=\mu+\epsilon_2$. At date~1, for $i>0$, LT1 compares selling now at $S_1-\eta$ with selling at date~2 at $S_2-\eta$, so its fee cancels; a market maker buys now at $S_1+\eta$ and expects to sell at $S_2-\eta$, a round trip:
\[ q_1^{LT1}=\frac{\mathbb{E}[S_2-\eta]-(S_1-\eta)}{\gamma\sigma^2}=\frac{\mu-S_1}{\gamma\sigma^2},\qquad
q_1^{MM}=\frac{\mathbb{E}[S_2-\eta]-(S_1+\eta)}{\gamma\sigma^2}=\frac{\mu-S_1-2\eta}{\gamma\sigma^2}. \]
Clearing, $i=nq_1^{MM}+q_1^{LT1}$, gives $(n+1)(\mu-S_1)-2n\eta=\gamma\sigma^2 i$, so
\[ S_1=\mu-\gamma\sigma^2\,\frac{i}{n+1}-\frac{2n}{n+1}\,\eta,\qquad
q_1^{LT1}=\frac{i}{n+1}+\frac{2n}{n+1}\,\frac{\eta}{\gamma\sigma^2}. \]
The fee deepens the discount and LT1 keeps more inventory than without fees: it pays its own fee directly and most of the market makers' round-trip fee through a lower price. If $\abs{i}$ is small relative to $\eta$ the desired trade can fall into a no-trade region.

\subsection*{Liquidity measures}
\emph{Price-impact slope.} Without fees, LT1's date-1 sale is $Q_1^{LT1}=i-q_1^{LT1}=\frac{n}{n+1}\,i$, so $\frac{i}{n+1}=\frac{Q_1^{LT1}}{n}$ and
\[ S_1=\mu-\frac{\gamma\sigma^2}{n}\,Q_1^{LT1}=\mu+\lambda Q_1^{LT1},\qquad \lambda=-\frac{\gamma\sigma^2}{n}, \]
with $Q_1^{LT1}$ the quantity sold. A more liquid market has a smaller $\abs{\lambda}$: more market makers, lower risk aversion or lower interim volatility.

\emph{Return reversal.} Add a date~0 and public news $\epsilon_1$ before LT1 trades, with efficient prices $\mu_0=\mathbb{E}[S_3]$, $\mu_1=\mathbb{E}[S_3\mid\epsilon_1]$, $\mu_2=\mathbb{E}[S_3\mid\epsilon_1,\epsilon_2]$. With $S_0=\mu_0$, $S_1=\mu_1+\lambda Q_1^{LT1}$ and $S_2=\mu_2$,
\[ \Delta_1=S_1-S_0=\epsilon_1+\lambda Q_1^{LT1},\qquad \Delta_2=S_2-S_1=\epsilon_2-\lambda Q_1^{LT1}. \]
If liquidity demand is independent of the news, all cross terms vanish except the order-pressure term:
\[ \operatorname{Cov}(\Delta_1,\Delta_2)=\operatorname{Cov}\big(\epsilon_1+\lambda Q_1^{LT1},\,\epsilon_2-\lambda Q_1^{LT1}\big)=-\lambda^2\operatorname{Var}(Q_1^{LT1})<0. \]
The first move contains order pressure and the second reverses it; as $\abs{\lambda}\to0$ the price process tends to the efficient-price martingale.

\section*{Kyle single-auction model}

Sources: \citep[chap.~``Strategic Trader Models I: Informed Traders'', pp.~91--98, eqs.~(4.3)--(4.5), (4.17), (4.20)]{catalogue-0496}; the setting as in \citep[chap.~``A Primer on the Microstructure of Financial Markets'', pp.~31--32]{crypto-cartea}, whose printed $\lambda$ on p.~33 is the reciprocal (corrected in its errata). Entry: \texttt{mod:kyle-single-auction}.

\subsection*{Assumptions}
\begin{enumerate}[label=\roman*.]
\setlength{\itemsep}{0pt}
\item The asset is worth $v\sim N(\mu,\sigma_v^2)$.
\item One risk-neutral insider observes $v$ and submits $x(v)$.
\item Price-insensitive noise traders submit $u\sim N(0,\sigma_u^2)$, independent of $v$; the insider does not see $u$.
\item Competitive risk-neutral market makers see only $y=x(v)+u$ and set $S(y)=\mathbb{E}[v\mid y]$, earning zero expected profit given $y$.
\item One auction, after which $v$ is realised. Equilibrium: the insider best-responds to the price rule, and the price rule is Bayes-consistent with the insider's strategy.
\end{enumerate}
The insider's profit is $\Pi=x(v)\big(v-S(x(v)+u)\big)$. Conjecture a linear equilibrium
\[ S(y)=\mu+\lambda y,\qquad x(v)=\beta(v-\mu). \]

\subsection*{The insider's best response}
Given the price rule and $\mathbb{E}[u]=0$,
\[ \mathbb{E}_u\big[x\big(v-\mu-\lambda(x+u)\big)\mid v\big]=x(v-\mu)-\lambda x^2, \]
strictly concave in $x$ for $\lambda>0$. The first-order condition $v-\mu-2\lambda x=0$ gives
\[ x^\star(v)=\frac{v-\mu}{2\lambda},\qquad \beta=\frac{1}{2\lambda}. \]

\subsection*{The market makers' pricing}
Under $x(v)=\beta(v-\mu)$, $y=\beta(v-\mu)+u$, and $(v,y)$ is jointly normal, so the conditional expectation is the linear projection
\[ \mathbb{E}[v\mid y]=\mu+\frac{\operatorname{Cov}(v,y)}{\operatorname{Var}(y)}\,y=\mu+\frac{\beta\sigma_v^2}{\beta^2\sigma_v^2+\sigma_u^2}\,y,
\qquad\text{so}\qquad \lambda=\frac{\beta\sigma_v^2}{\beta^2\sigma_v^2+\sigma_u^2}. \]

\subsection*{The fixed point}
Substitute $\beta=1/(2\lambda)$: $\lambda(\beta^2\sigma_v^2+\sigma_u^2)=\beta\sigma_v^2$ becomes $\frac{\sigma_v^2}{4\lambda}+\lambda\sigma_u^2=\frac{\sigma_v^2}{2\lambda}$, so $\lambda^2\sigma_u^2=\frac{\sigma_v^2}{4}$ and, taking the positive root,
\[ \lambda^\star=\frac{\sigma_v}{2\sigma_u},\qquad \beta^\star=\frac{\sigma_u}{\sigma_v},\qquad
x^\star(v)=\frac{\sigma_u}{\sigma_v}(v-\mu),\qquad S^\star(y)=\mu+\frac{\sigma_v}{2\sigma_u}\,y. \]
There is no other linear equilibrium; nonlinear equilibria are not excluded \citep[p.~97]{catalogue-0496}.

\subsection*{Profit, efficiency and revealed information}
Conditional on $v$, $\mathbb{E}_u[\Pi\mid v]=x^\star(v)(v-\mu)-\lambda^\star x^\star(v)^2=\frac{(v-\mu)^2}{2\lambda^\star}-\frac{(v-\mu)^2}{4\lambda^\star}=\frac{(v-\mu)^2}{4\lambda^\star}$, and ex ante
\[ \mathbb{E}[\Pi]=\frac{\sigma_v^2}{4\lambda^\star}=\frac{\sigma_v\sigma_u}{2}. \]
Market makers break even given $y$, $\mathbb{E}[v-S^\star(y)\mid y]=0$, so the price is semi-strong efficient for their information. In equilibrium $\beta^{\star2}\sigma_v^2=\sigma_u^2$, so $\operatorname{Var}(y)=2\sigma_u^2$ and $\operatorname{Cov}(v,y)=\sigma_u\sigma_v$:
\[ \operatorname{Var}(v\mid y)=\sigma_v^2-\frac{\operatorname{Cov}(v,y)^2}{\operatorname{Var}(y)}=\sigma_v^2-\frac{\sigma_u^2\sigma_v^2}{2\sigma_u^2}=\tfrac12\sigma_v^2, \]
whatever the order flow: one auction reveals half of the insider's information. The insider trades only part of the mispricing because a larger order reveals more through $y$; higher $\sigma_u$ lowers $\lambda^\star$ and raises the insider's order and profit, higher $\sigma_v$ raises $\lambda^\star$. The impact is informational and permanent, unlike the Grossman--Miller inventory pressure, which reverses.

\section*{Static limit-order depth}

Source: \citep[chap.~``A Primer on the Microstructure of Financial Markets'', pp.~28--30, eq.~(2.3)]{crypto-cartea}. Entry: \texttt{mod:static-limit-order-depth}.

A risk-neutral market maker posts a sell order at $S_t+\delta^+$ and a buy order at $S_t-\delta^-$. A buy (sell) market order arrives with probability $p^+$ ($p^-$) and reaches depth $\delta$ with probability $e^{-\kappa^\pm\delta}$; there is no adverse selection, inventory cost or reaction by other traders. Expected profit per decision is
\[ \mathbb{E}[\Pi(\delta^+,\delta^-)]=p^+e^{-\kappa^+\delta^+}\delta^+ + p^-e^{-\kappa^-\delta^-}\delta^-, \]
which separates by side. On each side, $\frac{d}{d\delta}\big(p\,e^{-\kappa\delta}\delta\big)=p\,e^{-\kappa\delta}(1-\kappa\delta)$, which is positive for $\delta<1/\kappa$ and negative beyond (own check: the second derivative at $\delta=1/\kappa$ is $-p\kappa e^{-1}<0$), so
\[ \delta^{\pm,\star}=\frac{1}{\kappa^\pm}, \]
the mean depth reached by market orders on that side. The arrival probabilities scale profit but not the optimal depth.

\section*{Roll's serial-covariance spread estimator}

Sources: \citep[chap.~``Volatility'', box, p.~415]{alpha-harris}; \citep[sec.~3.4, pp.~29--30]{s6-hasbrouck}. Entry: \texttt{res:roll-spread-estimator}.

Let $P_t=V_t+\tfrac12SQ_t$ with $V_t=V_{t-1}+\varepsilon_t$, $\mathbb{E}\varepsilon_t=0$, $\operatorname{Var}\varepsilon_t=\sigma^2$, and $Q_t=\pm1$ with probability one-half each, independent of earlier signs and of the innovations. Then
\[ \Delta P_t=\varepsilon_t+\tfrac12S(Q_t-Q_{t-1}). \]
\emph{Variance.} $\varepsilon_t$ is independent of the sign terms, and $\mathbb{E}[(Q_t-Q_{t-1})^2]=\mathbb{E}Q_t^2-2\mathbb{E}[Q_tQ_{t-1}]+\mathbb{E}Q_{t-1}^2=1-0+1=2$, so
\[ \operatorname{Var}(\Delta P_t)=\sigma^2+\tfrac14S^2\cdot2=\sigma^2+\tfrac12S^2. \]
\emph{First autocovariance.} The innovations are serially independent and independent of the signs, so only the sign terms contribute:
\[ \operatorname{Cov}(\Delta P_t,\Delta P_{t-1})=\tfrac14S^2\,\mathbb{E}\big[(Q_t-Q_{t-1})(Q_{t-1}-Q_{t-2})\big]
=\tfrac14S^2\big(\mathbb{E}[Q_tQ_{t-1}]-\mathbb{E}[Q_tQ_{t-2}]-\mathbb{E}Q_{t-1}^2+\mathbb{E}[Q_{t-1}Q_{t-2}]\big)=-\tfrac14S^2, \]
since every cross-product has zero expectation and $\mathbb{E}Q_{t-1}^2=1$. Autocovariances of order two or more are zero, because $\Delta P_t$ and $\Delta P_{t-k}$ share no sign or innovation for $k\ge2$. Hence, when the first autocovariance is negative,
\[ S=2\sqrt{-\operatorname{Cov}(\Delta P_t,\Delta P_{t-1})}. \]
How correlated signs and informed trading bias the estimate is in the entry's Pitfalls.

\section*{Source outside the catalogue}

\qfsource[https://www.nber.org/papers/w2641]{grossman-miller-1988-nber}{Sanford J. Grossman and Merton H. Miller, ``Liquidity and Market Structure'', NBER Working Paper 2641, July 1988; the working-paper version of \nolinkurl{grossman-miller-1988}, whose pages and numbering may differ.}

Every other cited key links to its entry in the pinned reading-list catalogue in the guide copy under \texttt{guide/}.

\end{document}
